% Ad Familiares 14.7 --- headnote
% ad-familiares-14-07, 7 June 49 BC, aboard ship in the harbour at Caieta

\textit{Cicero to his wife Terentia and daughter Tullia,
written aboard ship in the harbour at Caieta on the
seventh day before the Ides of June 49 BC (the
manuscript dateline: \textit{Scr. in portu Caietano
nave conscensa vii Id. Iun. a. 705 (49)}). After more
than four months of paralysed indecision in his Formian
villa --- weighing his consular dignity, his obligation
to Pompey, his promises to Caesar, his terror for his
family --- Cicero has at last embarked. The ship is in
the harbour; the line is about to be cast off; he is
sailing to join Pompey's camp in Epirus. He will not
see Italy again for almost a year, and the war he is
sailing into will be lost.}

\textit{The letter is short, and almost giddy with the
relief of decision. The four months of agony, he
declares, were bile: in the night he brought up
\textit{chol\=en akraton}, ``pure bile,'' and the moment
the body purged itself the soul cleared too. Give thanks
to Apollo and Aesculapius. The ship is good; he will
write to friends from on board commending Terentia and
Tullia to them; he need not urge them to courage,
knowing them braver than any man; he expects in time to
defend the commonwealth alongside men of his own kind.
The closing instructions are practical to the point of
domesticity --- mind your health, withdraw to country
houses well away from troops, the Arpinum estate will
carry the town household if grain prices rise, our
little Cicero sends his love. Then a doubled farewell,
and the date. He has not yet sailed; he is writing
between the decision and the wind.}
